\chapter{Problem Formulation}
\label{sec:problem-formulation}

\section{Proof States and Proof Scaffolds}

Given a theorem statement \(T\), the objective is to construct Lean code that proves \(T\) without using \texttt{sorry}, \texttt{admit}, or any temporary placeholder. Lean verification is the final correctness criterion: a proof is accepted only when the complete theorem is checked successfully.

GammaZero still uses proof states as the basic units of search. A Lean proof state provides the local context and target goal that the model must reason about. These remain part of the state representation. However, local context and goal alone are not enough for GammaZero, because they do not fully record where the goal sits inside the parent proof or how a completed child proof will be inserted back. This distinction is important because the same goal may occur in different parts of a proof and may have access to different local facts.

To preserve this missing structure, GammaZero adds a proof scaffold around each proof state. A scaffold is a partial Lean proof that contains the current hole to be solved, the surrounding proof code, and the structure that connects this hole back to the parent goal. The model is therefore not asked to solve a child goal as an independent theorem. Instead, it solves the goal inside the proof structure where the goal was created.

Each open state is assigned one target hole. The task of the state is to replace that hole with valid Lean code. Other unfinished holes in the same scaffold are treated as sibling holes. They are not solved by the current state; they belong to other states in the graph. This gives each state a clear local task while still keeping the full parent proof structure visible.

This design is especially important for skeleton actions. A skeleton action may introduce several child subgoals, but it also provides a partial proof structure showing how those subgoals are intended to support the parent goal. Solving the children separately is not enough by itself. Their proofs must later fit back into the scaffold so that the whole parent proof can be checked by Lean.

At a high level, a scaffolded proof state can be viewed as
\[
s = (C, i, \Gamma, G),
\]
where \(C\) is the proof scaffold, \(i\) identifies the target hole inside that scaffold, \(\Gamma\) is the local context, and \(G\) is the target goal. Lean provides \(\Gamma\) and \(G\) at the target hole, and these are shown to the model to make the local task explicit. However, the state is not represented only by \((\Gamma, G)\). The surrounding scaffold and target position are also part of the state because they record where the hole occurs and how the completed proof will connect back to the parent proof.
This distinction matters because two holes may have the same visible goal but occur in different scaffolds. A proof that is valid for one hole is not necessarily valid for the other.
This gives each state a simple rule: it may replace the target hole with Lean code, but it must leave the surrounding scaffold unchanged.

After the required child states are solved, GammaZero inserts their proof bodies back into the parent scaffold and asks Lean to verify the resulting proof again. Temporary placeholders used during search do not count as proof. They are only a way to check one part of a larger proof while the other parts are still being solved. Final success always requires complete Lean verification without placeholders.

\section{Action Space: Local Proof Actions and Skeleton Actions}

GammaZero separates two ways of continuing proof search. The first way is to solve the current proof state directly. The second way is to introduce proof structure that may create smaller subgoals. This gives two action types:
\[
\mathcal{A} = \mathcal{A}_{local} \cup \mathcal{A}_{skel}.
\]

\paragraph{Local proof actions.}
A local proof action is a generated Lean proof body for the current target hole. Its goal is to close the current proof state inside the current scaffold. The action may contain multiple tactics and may temporarily create subgoals during execution. However, it is accepted as a successful local proof action only if all of these goals are solved by the end of the generated proof body.

If a local proof action leaves any unresolved goal, it fails as a local proof action. It is not converted into a decomposition and it does not create child proof states in the search graph. This rule keeps the meaning of local proof actions simple: they either solve the current state or fail at that state.

This design reduces the number of search decisions inside one proof state. In a tactic-level search, each small tactic step may become a new search decision. In GammaZero, a local proof action represents a complete local attempt. Since it does not expand the graph, the growth of the graph is controlled mainly by skeleton actions.

\paragraph{Skeleton actions.}
A skeleton action is a generated Lean proof body that introduces proof structure with possible child subgoals. It may use constructs such as \texttt{have}, \texttt{suffices}, \texttt{cases}, or \texttt{induction}. Unlike a local proof action, a skeleton action is allowed to create child subgoals.

A skeleton action does not merely list subgoals. It also describes how the resulting proof pieces are intended to support the current goal. For example:
\begin{verbatim}
intro hp
have hq : Q := by
  sorry
exact h2 hq
\end{verbatim}
The statement \(Q\) becomes a child subgoal, while the last line shows how the completed proof of \(Q\) will be used to continue the parent proof.

A skeleton action may create zero or more child subgoals. If it creates no child subgoal and the current state is solved, it is simply a successful action. If it creates child subgoals, then all of them must eventually be solved for the skeleton to succeed. In the search graph, the child subgoals become child proof states.

This gives the graph an AND/OR structure. OR nodes are proof states, because a proof state can be solved by any successful outgoing action. Skeleton actions have AND semantics, because every child proof state created by the skeleton must be solved. This structure matches proof decomposition: a decomposition is useful only when all required subgoals can be completed and connected back to the parent proof.

The concrete running example for these definitions is moved to Appendix~\ref{app:running-example}.

\section{AND/OR Graph Semantics}

The action definitions above induce an AND/OR graph over scaffolded proof states. OR nodes are proof states: a state is solved when at least one outgoing action succeeds. Action nodes represent generated proof attempts. A local proof action is an action with no child proof state after success. A skeleton action may have child proof states, and it succeeds only when all required children are solved and the completed scaffold is accepted by Lean.

The graph makes two different relations explicit. A proof state represents a choice among candidate actions: any successful action can solve the state. A skeleton represents a decomposition: every child proof state created by the skeleton must be solved before the skeleton can solve its parent state.

GammaZero uses a graph rather than a tree because duplicate states can be shared. However, state identity is scaffold-sensitive. Two states are merged only when they refer to the same target problem inside the same proof scaffold. Thus, a child subgoal is not merged back into an ancestor merely because the visible goal text is the same. The scaffold records where the goal occurs and how it connects to the parent proof, so equal-looking goals in different scaffolds remain distinct.

Solved status propagates through this graph. When a child proof state is solved, its parent skeleton may become solved. When a skeleton action is solved, its parent proof state may become solved. This propagation can continue upward until the root theorem is solved.

This graph view is also used later for reward backup. OR nodes select among alternative actions, while skeleton actions are limited by their unfinished children. The full reward formulas are defined in Section~\ref{sec:reward}; here the important point is only the logical dependency: alternatives are OR choices, and decomposition children are AND requirements.

\section{Search Limits and Failure Labels}

GammaZero runs proof search under finite limits. These limits include the number of generated actions, the maximum search depth, the number of local proof attempts and skeleton attempts per state, the number of open states kept in the queue, model generation limits, and Lean verification timeouts. These limits are necessary because both model sampling and Lean verification are expensive.

Because the search is finite, the labels used in the graph are operational labels. They describe what happened during the current run; they do not describe mathematical truth. In particular, failing to solve a state within the current limits does not mean that the goal is false or impossible to prove.

The main status labels are:

\begin{center}
\begin{tabular}{ll}
\hline
\textbf{Status} & \textbf{Meaning} \\
\hline
\textsc{Solved} & The state or action has been closed by Lean without remaining placeholders. \\
\textsc{Open} & The state is still active and may be expanded later. \\
\textsc{Failed} & The action, state, or skeleton cannot continue successfully in the current run. \\
\textsc{Unresolved Within Limits} & The search stopped before the open state was solved. \\
\hline
\end{tabular}
\end{center}

A local proof action is marked successful only if Lean verifies that the current target has been fully solved. If it leaves an unresolved goal, it fails as a local proof action. The failed action may still be kept in the record for reward computation or analysis, but it does not create child proof states and does not solve the parent state.

A skeleton action has a stricter success condition. If it creates child proof states, then every child must be solved and the completed scaffold must be accepted by Lean. If one required child remains unsolved when the search gives up on that decomposition, the skeleton is marked failed in the graph. This does not mean that the child goal is mathematically impossible. It only means that the skeleton was not completed under the current search limits.

A proof state is solved when at least one outgoing action succeeds. A proof state may remain open while there are still allowed attempts or while one of its active decompositions is still being explored. If all useful attempts have failed or the state has exhausted its allowed search budget, it may be marked failed for the current run. If the global search stops while the state is still open, it is recorded as unresolved within limits.

This distinction is important for reward assignment. A repaired candidate may provide useful syntactic reward signal, but it cannot solve a state if placeholders remain. A skeleton with one solved child and one unresolved child may contain useful structure, but it is not a solved decomposition. The system may keep syntactic rewards and diagnostic information from such actions, while withholding delayed structural credit until the required proof states are actually solved.

Timeouts and system failures are treated in the same operational way. They are recorded as failed actions or interrupted checks for the current run, not as evidence about the truth of the theorem. The final correctness boundary remains Lean verification of the completed proof: the root theorem is solved only when Lean accepts the fully assembled proof without \texttt{sorry}, \texttt{admit}, or any temporary placeholder.

\section{Design Invariants}

The formulation is governed by a small set of invariants. First, every scaffolded proof state has exactly one target hole. This target is the only location that the current state is responsible for solving. Other unfinished holes in the same scaffold belong to sibling states.

Second, while one target is being solved, sibling subgoals may be represented by temporary \texttt{admit} placeholders. These placeholders are part of the surrounding scaffold and must not be changed by the candidate action. The current action may replace only the target hole.

Third, local proof actions are terminal with respect to the search graph. They either close the current target or remain failed attempts recorded at the same state; they never introduce child proof states.

Fourth, a skeleton action may create child proof states by extending the current scaffold with child subgoals. If it creates child states, the skeleton cannot solve its parent until every required child has been solved and the completed scaffold passes Lean verification. If a skeleton creates no child state, it is solved only when Lean verifies that it directly closes the current target.

Fifth, a proof state can be marked \textsc{Solved} only through Lean verification. For a local state, the replacement for the target hole must not contain \texttt{sorry}, \texttt{admit}, or any temporary placeholder. For a parent skeleton, all child proof bodies must be stitched back into the parent scaffold and the completed block must be checked again.

The global invariant is that temporary placeholders are never part of an accepted theorem. They are used only to isolate responsibilities during search. A repaired script that still contains placeholders may provide reward information, but it cannot solve a state or appear in the final proof. The root theorem is accepted only after recursive assembly reaches the root and Lean verifies the complete theorem without unresolved placeholders.

These invariants also clarify how failure is represented. Parser errors, elaboration errors, timeouts, and policy-check rejections are failed actions, not failed theorems. An open state that is not solved within the current search limits remains open from the mathematical point of view, even if the exported record labels it as \textsc{Unresolved Within Limits}. This distinction is important for future training: the model should learn that the explored branch did not succeed under the current search configuration, but it should not learn that the underlying goal is false.
